\documentclass[a4paper,oneside,11pt]{book}

% ===== Package =====
\usepackage[T1]{fontenc}
\usepackage[utf8]{inputenc}
\usepackage{lmodern}
\usepackage[indonesian]{babel}
\usepackage[htt]{hyphenat}
\usepackage{microtype}
\usepackage{graphicx}
\usepackage{titling}
\usepackage{geometry}
\usepackage{listings}
\usepackage{xcolor}
\usepackage{float}
\usepackage{array}
\usepackage{tabularx}
\usepackage{caption}
\usepackage{indentfirst}
\usepackage{xurl}
\usepackage[hidelinks]{hyperref}

% ===== Margin =====
\geometry{
    left=3cm,
    right=2.5cm,
    top=3cm,
    bottom=3cm
}

% ===== Pengaturan teks =====
\emergencystretch=3em
\tolerance=1500
\hbadness=5000
\linespread{1.15}
\setlength{\parskip}{3pt}
\renewcommand{\_}{\textunderscore\allowbreak{}}
\captionsetup{font=small,labelfont=bf}

% ===== Pengaturan kode =====
\definecolor{codegreen}{rgb}{0,0.6,0}
\definecolor{codegray}{rgb}{0.5,0.5,0.5}
\definecolor{codepurple}{rgb}{0.58,0,0.82}
\definecolor{backcolour}{rgb}{0.95,0.95,0.92}

\lstdefinestyle{codestyle}{
    backgroundcolor=\color{backcolour},
    commentstyle=\color{codegreen},
    keywordstyle=\color{codepurple},
    numberstyle=\tiny\color{codegray},
    stringstyle=\color{codegreen},
    basicstyle=\ttfamily\footnotesize,
    breakatwhitespace=false,
    breaklines=true,
    postbreak=\mbox{\textcolor{codegray}{$\hookrightarrow$}\space},
    captionpos=b,
    keepspaces=true,
    numbers=left,
    numbersep=6pt,
    xleftmargin=2em,
    framexleftmargin=1.6em,
    xrightmargin=0.3em,
    showspaces=false,
    showstringspaces=false,
    showtabs=false,
    tabsize=2,
    frame=single,
    aboveskip=1em,
    belowskip=0.8em
}

\lstset{style=codestyle}
\AtBeginDocument{\renewcommand{\lstlistingname}{Kode}}

% ===== Identitas =====
\title{
    Laporan Praktikum\\
    Praktikum Pemrograman Web 2\\
    Pertemuan ke-8\\
    Seeder dan Pagination pada Aplikasi CRUD Project
}

\author{
    Nama Mahasiswa: Gayuh Mukti Aji\\
    NIM: 25/566296/SV/27058
}

\date{27 Oktober 2026}

\begin{document}

% ===== Halaman judul =====
\begin{titlingpage}
    \begin{center}
        \vspace*{1cm}

        {\LARGE\bfseries \thetitle\par}

        \vspace{1.5cm}

        % Letakkan lambang_ugm.png pada folder yang sama.
        \IfFileExists{lambang_ugm.png}{
            \includegraphics[height=6cm]{lambang_ugm.png}
        }{
            \fbox{
                \parbox[c][6cm][c]{6cm}{
                    \centering
                    Lambang Universitas Gadjah Mada
                }
            }
        }

        \vspace{1.5cm}

        {\large
            \theauthor\\
            Kelas PL3B2\par
        }

        \vspace{0.8cm}

        % Isi sesuai dosen pengampu mata kuliah.
        Dosen Pengampu: \underline{\hspace{5cm}}

        \vspace{1cm}

        \thedate
    \end{center}
\end{titlingpage}

\tableofcontents
\newpage

% =========================================================
\chapter{Tujuan Praktikum}

Tujuan praktikum pertemuan ke-8 adalah sebagai berikut:

\begin{enumerate}
    \item Memahami fungsi \textit{seeder} untuk mengisi basis data dengan data contoh.
    \item Membuat dan menjalankan seeder secara terpisah maupun melalui \texttt{DatabaseSeeder}.
    \item Menerapkan pagination pada daftar project menggunakan \texttt{paginate()}.
    \item Menampilkan informasi halaman, jumlah data, dan tautan navigasi.
    \item Menggunakan tampilan pagination Bootstrap.
    \item Membuat nomor urut yang berlanjut antarhalaman.
    \item Menambahkan pencarian judul yang tetap aktif ketika pengguna berpindah halaman.
    \item Menyediakan pilihan jumlah data sebanyak 5, 10, atau 25 per halaman.
\end{enumerate}

% =========================================================
\chapter{Dasar Teori}

\section{Seeder}

Seeder merupakan kelas untuk mengisi basis data dengan data contoh.
Pada Laravel, kelas seeder berada di direktori
\path{database/seeders/}.
Seeder membantu proses pengembangan karena data tidak perlu dimasukkan
satu per satu melalui form \cite{laravel-seeding}.

Seeder berbeda dengan migration. Migration digunakan untuk mengatur
struktur tabel, sedangkan seeder digunakan untuk mengisi data ke dalam
tabel yang sudah tersedia \cite{laravel-seeding,laravel-migrations}.

\section{DatabaseSeeder}

\texttt{DatabaseSeeder} menjadi tempat untuk memanggil beberapa seeder.
Pemanggilan dilakukan melalui method \texttt{\$this->call()}.
Dengan cara ini, beberapa seeder dapat dijalankan melalui satu perintah \cite{laravel-seeding}:

\begin{lstlisting}[language={},caption={Menjalankan DatabaseSeeder}]
php artisan db:seed
\end{lstlisting}

Seeder tertentu dapat dijalankan menggunakan pilihan
\texttt{--class}, misalnya:

\begin{lstlisting}[language={},caption={Menjalankan satu seeder}]
php artisan db:seed --class=ProjectTableSeeder
\end{lstlisting}

\section{Pencegahan Data Ganda}

Pemanggilan \texttt{create()} akan membuat baris baru setiap kali seeder
dijalankan. Apabila tidak ditangani, seeder yang dijalankan berulang
dapat menambah data contoh yang sama.

Pada praktikum ini digunakan \texttt{firstOrCreate()}.
Method tersebut mencari data berdasarkan atribut yang diberikan.
Apabila data ditemukan, data yang ada digunakan kembali. Apabila belum
ditemukan, Laravel membuat data baru \cite{laravel-eloquent}.

Pencarian data contoh menggunakan judul project. Dengan demikian,
seeder dapat dijalankan ulang tanpa mengosongkan tabel.

\section{Faker}

Faker digunakan untuk menghasilkan data contoh, seperti kalimat dan
paragraf. Pada praktikum ini, Faker menghasilkan judul dan deskripsi
untuk seeder kedua.

Generator menggunakan nilai awal tetap melalui \texttt{seed(8)}.
Tujuannya agar urutan data acak yang dihasilkan tetap konsisten ketika
seeder dijalankan ulang pada lingkungan dan versi Faker yang sama \cite{faker}.
Judul yang konsisten dapat dicari kembali oleh
\texttt{firstOrCreate()} \cite{laravel-eloquent}.

\section{Pagination}

Pagination membagi daftar data menjadi beberapa halaman.
Sebagai contoh, 35 data dengan batas 10 data per halaman akan
ditampilkan dalam empat halaman.

Laravel menyediakan \texttt{paginate()} untuk mengambil data sesuai
halaman yang diminta. Nomor halaman dibaca dari parameter
\texttt{page} pada URL, misalnya \path{/projects?page=2} \cite{laravel-pagination}.

Hasil \texttt{paginate()} berupa objek
\texttt{LengthAwarePaginator}. Objek ini menyediakan data halaman aktif
beserta informasi jumlah seluruh hasil.

\section{Informasi dan Navigasi Halaman}

Beberapa method yang digunakan dalam pagination adalah:

\begin{itemize}
    \item \texttt{currentPage()}: memperoleh nomor halaman aktif.
    \item \texttt{lastPage()}: memperoleh nomor halaman terakhir.
    \item \texttt{total()}: memperoleh jumlah seluruh data hasil query.
    \item \texttt{perPage()}: memperoleh batas data per halaman.
    \item \texttt{firstItem()}: memperoleh nomor urut data pertama pada halaman aktif.
    \item \texttt{links()}: menampilkan tautan navigasi halaman.
\end{itemize}

Ketika pencarian diterapkan, nilai \texttt{total()} menunjukkan jumlah
hasil pencarian, bukan jumlah seluruh project tanpa filter \cite{laravel-pagination}.

\section{Parameter Pencarian}

Pencarian judul menggunakan kondisi \texttt{LIKE} pada query.
Karakter \texttt{\%} digunakan agar kata pencarian dapat ditemukan
pada bagian judul \cite{laravel-queries}.

Method \texttt{withQueryString()} mempertahankan parameter query pada
tautan pagination. Dengan demikian, parameter pencarian dan jumlah
data per halaman tetap digunakan ketika pengguna membuka halaman
berikutnya \cite{laravel-pagination}.

% =========================================================
\chapter{Hasil dan Pembahasan}

\section{Gambaran Proyek}

Praktikum diterapkan pada proyek Laravel yang sebelumnya telah memiliki
fitur CRUD project. Data project disimpan pada tabel
\texttt{projects} dengan atribut utama \texttt{title} dan
\texttt{description}.

Modul menggunakan contoh data \texttt{Post}, sedangkan proyek ini
menggunakan \texttt{Project}. Penyesuaian dilakukan pada nama model,
seeder, controller, dan view tanpa mengubah tujuan praktikum.

Tampilan tetap menggunakan Bootstrap sederhana. Daftar project
ditampilkan dalam tabel, sedangkan pencarian dan pilihan jumlah data
ditempatkan di atas tabel.

\begin{table}[H]
    \centering
    \caption{Komponen implementasi pertemuan ke-8}
    \label{tab:komponen}
    \renewcommand{\arraystretch}{1.3}
    \begin{tabularx}{\textwidth}{
        |>{\raggedright\arraybackslash}p{5cm}
        |>{\raggedright\arraybackslash}X|
    }
        \hline
        \textbf{Komponen} & \textbf{Fungsi} \\ \hline
        \texttt{ProjectTableSeeder}
        & Membuat 10 project contoh dengan judul berurutan. \\ \hline

        \texttt{RandomProject\-TableSeeder}
        & Membuat 25 project dengan judul acak. \\ \hline

        \texttt{DatabaseSeeder}
        & Memanggil kedua seeder project. \\ \hline

        \texttt{ProjectController}
        & Mengatur pencarian, pengurutan, dan pagination. \\ \hline

        \texttt{projects/index.blade.php}
        & Menampilkan tabel, form pencarian, informasi halaman, dan navigasi. \\ \hline

        \texttt{AppServiceProvider}
        & Mengaktifkan tampilan pagination Bootstrap 5. \\ \hline
    \end{tabularx}
\end{table}

\section{Model Project}

Model \texttt{Project} sudah memiliki properti
\texttt{\$fillable}. Properti ini menentukan atribut yang dapat diisi
melalui operasi seperti \texttt{create()} dan
\texttt{firstOrCreate()} \cite{laravel-eloquent}.

\begin{lstlisting}[language=PHP,caption={Model Project}]
<?php

namespace App\Models;

use Illuminate\Database\Eloquent\Model;

class Project extends Model
{
    protected $fillable = ['title', 'description'];
}
\end{lstlisting}

Atribut yang digunakan seeder adalah \texttt{title} dan
\texttt{description}, sehingga keduanya dicantumkan pada
\texttt{\$fillable}.

\section{Membuat Seeder untuk 10 Project}

Seeder pertama berada pada
\path{database/seeders/ProjectTableSeeder.php}.
Data dibuat menggunakan perulangan dari 1 sampai 10.

\begin{lstlisting}[language=PHP,caption={ProjectTableSeeder.php}]
<?php

namespace Database\Seeders;

use App\Models\Project;
use Illuminate\Database\Seeder;

class ProjectTableSeeder extends Seeder
{
    public function run(): void
    {
        for ($i = 1; $i <= 10; $i++) {
            Project::firstOrCreate(
                ['title' => 'Project Judul Ke-'.$i],
                ['description' => 'Deskripsi untuk project judul ke-'.$i]
            );
        }
    }
}
\end{lstlisting}

Argumen pertama \texttt{firstOrCreate()} menjadi kondisi pencarian,
yaitu judul project. Argumen kedua berisi nilai tambahan untuk data
baru.

Sebagai contoh, iterasi pertama menghasilkan judul
\texttt{Project Judul Ke-1}. Jika judul tersebut sudah tersedia,
seeder menggunakan data yang ada dan tidak menambah baris baru.

\section{Membuat Seeder untuk 25 Project Acak}

Latihan pada modul dikembangkan melalui seeder kedua, yaitu
\path{database/seeders/RandomProjectTableSeeder.php}.

\begin{lstlisting}[language=PHP,caption={RandomProjectTableSeeder.php}]
<?php

namespace Database\Seeders;

use App\Models\Project;
use Faker\Factory;
use Illuminate\Database\Seeder;

class RandomProjectTableSeeder extends Seeder
{
    public function run(): void
    {
        $faker = Factory::create('id_ID');

        // Seed tetap membuat judul acak yang sama
        // saat seeder dijalankan ulang.
        $faker->seed(8);

        for ($i = 1; $i <= 25; $i++) {
            Project::firstOrCreate(
                ['title' => 'Project '.$faker->unique()->sentence(3)],
                ['description' => $faker->paragraph()]
            );
        }
    }
}
\end{lstlisting}

Penjelasan kode tersebut adalah sebagai berikut:

\begin{enumerate}
    \item \texttt{Factory::create('id\_ID')} membuat generator Faker dengan locale Indonesia.
    \item \texttt{seed(8)} menetapkan nilai awal generator.
    \item Perulangan menghasilkan 25 data contoh.
    \item \texttt{unique()->sentence(3)} menghasilkan kalimat judul yang berbeda dalam pemanggilan generator tersebut.
    \item \texttt{paragraph()} menghasilkan deskripsi.
    \item \texttt{firstOrCreate()} mencegah penambahan ulang judul contoh yang sudah tersedia.
\end{enumerate}

Pada basis data kosong, kedua seeder menghasilkan 35 project.
Jika terdapat data manual sebelumnya, jumlah keseluruhan project dapat
lebih banyak karena tabel tidak dikosongkan.

\section{Memanggil Seeder melalui DatabaseSeeder}

Kedua seeder dipanggil melalui \texttt{DatabaseSeeder}.
Pembuatan pengguna contoh yang sebelumnya sudah ada juga diberi
pemeriksaan email agar tidak dibuat kembali.

\begin{lstlisting}[language=PHP,caption={DatabaseSeeder.php}]
<?php

namespace Database\Seeders;

use App\Models\User;
use Illuminate\Database\Console\Seeds\WithoutModelEvents;
use Illuminate\Database\Seeder;

class DatabaseSeeder extends Seeder
{
    use WithoutModelEvents;

    public function run(): void
    {
        if (! User::where('email', 'test@example.com')->exists()) {
            User::factory()->create([
                'name' => 'Test User',
                'email' => 'test@example.com',
            ]);
        }

        $this->call([
            ProjectTableSeeder::class,
            RandomProjectTableSeeder::class,
        ]);
    }
}
\end{lstlisting}

Pada pelaksanaan praktikum, status migrasi diperiksa terlebih dahulu.
Setelah seluruh migrasi tercatat selesai, seeder dijalankan melalui
perintah berikut:

\begin{lstlisting}[language={},caption={Pemeriksaan migrasi dan pengisian data contoh}]
php artisan migrate:status
php artisan db:seed
\end{lstlisting}

Kedua seeder berhasil dijalankan pada MySQL lokal.
Pengisian data tidak menggunakan \texttt{truncate()} atau
\texttt{migrate:fresh}, sehingga data project yang sudah tersedia
tetap dipertahankan.

\section{Mengubah Daftar Project menjadi Pagination}

Sebelumnya, method \texttt{index()} menggunakan
\texttt{Project::all()}, sehingga seluruh project diambil sekaligus.
Pada pertemuan ini, pengambilan data diubah menggunakan
\texttt{paginate()}.

Berikut method \texttt{index()} pada
\path{app/Http/Controllers/ProjectController.php}.

\begin{lstlisting}[language=PHP,caption={Method index pada ProjectController}]
public function index(Request $request)
{
    $validated = $request->validate([
        'search' => 'nullable|string|max:200',
        'per_page' => 'nullable|integer|in:5,10,25',
    ]);

    $search = trim($validated['search'] ?? '');
    $perPage = (int) ($validated['per_page'] ?? 10);

    $query = Project::query();

    if ($search !== '') {
        $query->where('title', 'like', '%'.$search.'%');
    }

    $projects = $query->orderByDesc('created_at')
        ->orderByDesc('id')
        ->paginate($perPage)
        ->withQueryString();

    return view(
        'projects.index',
        compact('projects', 'search', 'perPage')
    );
}
\end{lstlisting}

Alur method tersebut adalah sebagai berikut:

\begin{enumerate}
    \item Memvalidasi parameter pencarian dan jumlah data per halaman.
    \item Menggunakan pencarian kosong dan batas 10 data sebagai nilai awal.
    \item Membentuk query melalui \texttt{Project::query()}.
    \item Menambahkan kondisi pencarian apabila pengguna mengisi judul.
    \item Mengurutkan data berdasarkan tanggal pembuatan terbaru.
    \item Menggunakan ID sebagai urutan tambahan apabila tanggal pembuatan sama.
    \item Mengambil data halaman aktif dengan \texttt{paginate()}.
    \item Mempertahankan parameter URL melalui \texttt{withQueryString()}.
    \item Mengirim hasil dan nilai filter ke view.
\end{enumerate}

Urutan tambahan berdasarkan ID diperlukan karena beberapa data seeder
dapat memiliki waktu pembuatan yang sama. Dengan urutan tersebut,
penampilan data pada pengujian menjadi konsisten.

\section{Form Pencarian dan Pilihan Jumlah Data}

Form pada halaman daftar menggunakan method \texttt{GET}.
Pengguna dapat mencari judul sekaligus menentukan jumlah data per
halaman.

Berikut cuplikan bagian utama form pada
\path{resources/views/projects/index.blade.php}.

\begin{lstlisting}[language=HTML,caption={Form pencarian dan pilihan jumlah data}]
<form action="{{ route('projects.index') }}"
    method="GET" class="row g-2 mb-3">

    <div class="col-sm-6">
        <label for="search" class="form-label">
            Cari judul
        </label>
        <input type="search" name="search" id="search"
            class="form-control"
            value="{{ old('search', $search) }}"
            maxlength="200">
    </div>

    <div class="col-sm-3">
        <label for="per_page" class="form-label">
            Data per halaman
        </label>
        <select name="per_page" id="per_page"
            class="form-select">
            @foreach ([5, 10, 25] as $size)
                <option value="{{ $size }}"
                    @selected($perPage == $size)>
                    {{ $size }}
                </option>
            @endforeach
        </select>
    </div>

    <div class="col-sm-3 d-flex align-items-end gap-2">
        <button type="submit" class="btn btn-primary">
            Cari
        </button>
        <a href="{{ route('projects.index') }}"
            class="btn btn-secondary">
            Reset
        </a>
    </div>
</form>
\end{lstlisting}

Sebagai contoh, pencarian kata \texttt{Judul} dengan batas lima data
menghasilkan parameter berikut:

\begin{lstlisting}[language={},caption={Contoh parameter pencarian}]
/projects?search=Judul&per_page=5
\end{lstlisting}

Ketika pengguna membuka halaman kedua, parameter pencarian tetap
disertakan:

\begin{lstlisting}[language={},caption={Contoh parameter pada halaman kedua}]
/projects?search=Judul&per_page=5&page=2
\end{lstlisting}

Tombol Reset membuka kembali daftar tanpa parameter pencarian.
Form menggunakan \texttt{GET} karena proses ini membaca data.

\section{Nomor Urut yang Berlanjut}

Nomor urut sebelumnya menggunakan \texttt{\$loop->iteration}.
Cara tersebut akan memulai penomoran dari 1 pada setiap halaman.
Untuk membuat nomor berlanjut, digunakan
\texttt{firstItem()} dan indeks perulangan.

\begin{lstlisting}[language=HTML,caption={Nomor urut pada tabel project}]
@forelse ($projects as $project)
    <tr>
        <td>
            {{ $projects->firstItem() + $loop->index }}
        </td>
        <td>{{ $project->title }}</td>
        <td>
            <a href="{{ route('projects.show', $project) }}">
                Lihat Detail
            </a>
        </td>
    </tr>
@empty
    <tr>
        <td colspan="3">
            @if ($projects->total() > 0)
                Tidak ada data pada halaman ini.
            @elseif ($search !== '')
                Tidak ada project yang cocok.
            @else
                Belum ada data project.
            @endif
        </td>
    </tr>
@endforelse
\end{lstlisting}

Pada halaman kedua dengan batas 10 data, nilai
\texttt{firstItem()} adalah 11.
Karena \texttt{\$loop->index} dimulai dari 0, nomor baris pertama
menjadi 11, kemudian 12, dan seterusnya.

Jika batas data diubah menjadi lima, nomor baris pertama halaman kedua
menjadi 6. Penomoran mengikuti hasil pagination yang sedang digunakan.

\section{Informasi dan Tautan Halaman}

Informasi halaman dan navigasi ditempatkan setelah tabel, masih di
dalam bagian \texttt{content}.

\begin{lstlisting}[language=HTML,caption={Informasi halaman dan tautan navigasi}]
<p>
    Halaman:
    {{ $projects->currentPage() }}
    dari {{ $projects->lastPage() }}<br>

    Jumlah Data: {{ $projects->total() }}<br>

    Data Per Halaman: {{ $projects->perPage() }}
</p>

{{ $projects->links() }}
\end{lstlisting}

Pada pengujian dengan 35 project dan batas 10 data, halaman pertama
menampilkan informasi:

\begin{lstlisting}[language={},caption={Informasi pada basis data pengujian}]
Halaman: 1 dari 4
Jumlah Data: 35
Data Per Halaman: 10
\end{lstlisting}

Apabila pencarian hanya menemukan 10 project, jumlah data dan jumlah
halaman dihitung berdasarkan 10 hasil tersebut.

\section{Menggunakan Pagination Bootstrap}

Pagination menggunakan tampilan Bootstrap 5 agar sesuai dengan
komponen tabel dan form pada aplikasi.
Pengaturan ditempatkan pada
\path{app/Providers/AppServiceProvider.php}.

\begin{lstlisting}[language=PHP,caption={Pengaturan pagination Bootstrap}]
use Illuminate\Pagination\Paginator;

// Di dalam kelas AppServiceProvider:
public function boot(): void
{
    Paginator::useBootstrapFive();
}
\end{lstlisting}

Setelah pengaturan tersebut, pemanggilan
\texttt{\$projects->links()} menggunakan view pagination Bootstrap.
Tidak diperlukan penyalinan seluruh view pagination ke dalam proyek \cite{laravel-pagination}.


% =========================================================
\chapter{Kesimpulan}

Berdasarkan praktikum yang telah dilakukan, seeder dapat digunakan
untuk menyediakan data contoh secara lebih cepat.
Pada proyek ini dibuat dua seeder, yaitu seeder untuk 10 project
berurutan dan seeder untuk 25 project dengan judul acak.
Keduanya dipanggil melalui \texttt{DatabaseSeeder} dan telah berhasil
dijalankan pada MySQL lokal.

Pagination diterapkan melalui \texttt{paginate()} pada daftar
project. Daftar menampilkan informasi halaman, jumlah hasil, batas
data, serta tautan navigasi dengan Bootstrap.
Nomor urut dibuat berlanjut menggunakan \texttt{firstItem()}
dan indeks perulangan.

Latihan dikembangkan melalui pencarian judul dan pilihan batas
5, 10, atau 25 data per halaman.
Parameter pencarian dipertahankan menggunakan
\texttt{withQueryString()}, sehingga pengguna dapat berpindah halaman
tanpa kehilangan filter.

Seluruh pengujian otomatis menghasilkan 26 pengujian lulus dengan
167 assertions. Proses build frontend juga berhasil.
Hasil tersebut menunjukkan bahwa implementasi CRUD sebelumnya tetap
berjalan bersama penambahan Seeder dan Pagination.

% =========================================================
% Daftar pustaka dibentuk otomatis oleh thebibliography.
\cleardoublepage
\renewcommand{\bibname}{Daftar Pustaka}

\begin{thebibliography}{9}
\addcontentsline{toc}{chapter}{Daftar Pustaka}
\raggedright
\setlength{\itemsep}{0.8em}

\bibitem{laravel-seeding}
Laravel, ``Database: Seeding,''
\textit{Laravel 12.x Documentation}. [Daring].
Tersedia: \url{https://laravel.com/docs/12.x/seeding}.
[Diakses: 27-Okt-2026].

\bibitem{laravel-migrations}
Laravel, ``Database: Migrations,''
\textit{Laravel 12.x Documentation}. [Daring].
Tersedia: \url{https://laravel.com/docs/12.x/migrations}.
[Diakses: 27-Okt-2026].

\bibitem{laravel-eloquent}
Laravel, ``Eloquent: Getting Started,''
\textit{Laravel 12.x Documentation}. [Daring].
Tersedia: \url{https://laravel.com/docs/12.x/eloquent}.
[Diakses: 27-Okt-2026].

\bibitem{faker}
FakerPHP, ``Faker,''
\textit{FakerPHP Documentation}. [Daring].
Tersedia: \url{https://fakerphp.org/}.
[Diakses: 27-Okt-2026].

\bibitem{laravel-pagination}
Laravel, ``Database: Pagination,''
\textit{Laravel 12.x Documentation}. [Daring].
Tersedia: \url{https://laravel.com/docs/12.x/pagination}.
[Diakses: 27-Okt-2026].

\bibitem{laravel-queries}
Laravel, ``Database: Query Builder,''
\textit{Laravel 12.x Documentation}. [Daring].
Tersedia: \url{https://laravel.com/docs/12.x/queries}.
[Diakses: 27-Okt-2026].

\end{thebibliography}

\end{document}